\subsection{Гибридный NEAT с самовниманием}


\subsubsection{Контекст и проблема}

Классический алгоритм NEAT, несмотря на свою эффективность в задачах управления с низкоразмерным входом, сталкивается с трудностями при масштабировании на задачи с высокой размерностью входных данных (например, визуальные входы или сложные сенсорные массивы). В таких условиях пространство поиска топологий становится слишком большим, что замедляет эволюцию и часто приводит к субоптимальным, переусложненным структурам.


\subsubsection{Предложенный метод}

В работе \cite{khamesian2023} предлагают модификацию Hybrid Self-Attention NEAT, которая интегрирует механизм самовнимания (Self-Attention), заимствованный из архитектур трансформеров, непосредственно в процесс нейроэволюции.​

Ключевые особенности подхода:

\begin{itemize}

  \item \textbf{Механизм внимания как непрямое кодирование:} Вместо того чтобы жестко эволюционировать связи для каждого пикселя или входа, алгоритм использует слой самовнимания для динамического взвешивания важности входных признаков. Это позволяет агенту фокусироваться на значимых частях состояния среды.

  \item \textbf{Гибридная оптимизация:} Работа объединяет эволюционный поиск архитектуры (NEAT) с элементами обучения, специфичными для механизмов внимания, что позволяет эффективно находить компактные представления для сложных входов.

\end{itemize}


\subsubsection{Результаты и выводы}

Метод был протестирован на классических бенчмарках (игры Atari), требующих обработки визуальной информации.

\begin{itemize}

  \item \textbf{Эффективность параметров:} Модель показала способность решать задачи, используя на порядки меньше параметров, чем традиционные методы глубокого обучения с подкреплением (Deep RL), такие как DQN. % todo сделать footnote?

  \item \textbf{Интерпретируемость:} Использование внимания повышает интерпретируемость действий агента, так как можно отследить, на какие части входных данных агент опирался при принятии решения.

\end{itemize}


\subsubsection{Связь с текущим исследованием}

Данная работа демонстрирует современное состояние (SOTA) развития нейроэволюционных алгоритмов. Она подтверждает, что NEAT не является устаревшим методом, а служит гибкой базой для построения сложных когнитивных архитектур. В рамках диплома эта статья обосновывает перспективность использования NEAT для моделирования когнитивных агентов, так как показывает потенциал метода к масштабированию и интеграции с современными механизмами внимания (attention mechanisms), свойственными человеческому восприятию.